% **************************************************************************************************************
% A Classic Thesis Style
% An Homage to The Elements of Typographic Style
%
% Copyright (C) 2011 Andr\'e Miede http://www.miede.de
%
% If you like the style then I would appreciate a postcard. My address 
% can be found in the file ClassicThesis.pdf. A collection of the 
% postcards I received so far is available online at 
% http://postcards.miede.de
%
% License:
% This program is free software; you can redistribute it and/or modify
% it under the terms of the GNU General Public License as published by
% the Free Software Foundation; either version 2 of the License, or
% (at your option) any later version.
%
% This program is distributed in the hope that it will be useful,
% but WITHOUT ANY WARRANTY; without even the implied warranty of
% MERCHANTABILITY or FITNESS FOR A PARTICULAR PURPOSE.  See the
% GNU General Public License for more details.
%
% You should have received a copy of the GNU General Public License
% along with this program; see the file COPYING.  If not, write to
% the Free Software Foundation, Inc., 59 Temple Place - Suite 330,
% Boston, MA 02111-1307, USA.
%
% **************************************************************************************************************
% Note:
%    * You must not use "u etc. in strings/commands that will be spaced out (use \"u or real umlauts instead)
%    * New enumeration (small caps): \begin{aenumerate} \end{aenumerate}
%    * For margin notes: \marginpar or \graffito{}
%    * Do not use bold fonts in this style, it is designed around them
%    * Use tables as in the examples
%    * See classicthesis-preamble.sty for useful commands
% **************************************************************************************************************
% To Do:
%    * Eliminare l'opzione draft e drafting dal comando \documentclass
%    * [Opzionale] Aggiungere le opzioni twoside + BCOR=5mm e "decommentare" i margini nella pagina di titolo
% **************************************************************************************************************
\documentclass[a4paper,openright,numbers=noenddot,headinclude,footinclude,abstracton,11pt]{scrreprt}
\usepackage[italian]{babel}
\usepackage{makeidx}
\usepackage{parskip}
\usepackage{textcomp}
\usepackage[utf8]{inputenc}
% ****************************************************************************************************
% classicthesis-config.tex 
% formerly known as loadpackages.sty, classicthesis-ldpkg.sty, and classicthesis-preamble.sty 
% Use it at the beginning of your ClassicThesis.tex, or as a LaTeX Preamble 
% in your ClassicThesis.{tex,lyx} with \input{classicthesis-config}
% ****************************************************************************************************  
% If you like the classicthesis, then I would appreciate a postcard. 
% My address can be found in the file ClassicThesis.pdf. A collection 
% of the postcards I received so far is available online at 
% http://postcards.miede.de
% ****************************************************************************************************

% ****************************************************************************************************
% 1. Configure classicthesis for your needs here, e.g., remove "drafting" below 
% in order to deactivate the time-stamp on the pages
% ****************************************************************************************************
\PassOptionsToPackage{eulerchapternumbers,listings,%drafting,%
				 pdfspacing,%floatperchapter,%linedheaders,%
				 subfig,beramono,eulermath,parts}{classicthesis}										
% ********************************************************************
% Available options for classicthesis.sty 
% (see ClassicThesis.pdf for more information):
% drafting
% parts nochapters linedheaders
% eulerchapternumbers beramono eulermath pdfspacing minionprospacing
% tocaligned dottedtoc manychapters
% listings floatperchapter subfig
% ********************************************************************

% ********************************************************************
% Triggers for this config
% ******************************************************************** 
\usepackage{ifthen}
\newboolean{enable-backrefs} % enable backrefs in the bibliography
\setboolean{enable-backrefs}{false} % true false
% ****************************************************************************************************


% ****************************************************************************************************
% 2. Personal data and user ad-hoc commands
% ****************************************************************************************************
\newcommand{\myTitle}{Relazione di stage\xspace}
\newcommand{\mySubtitle}{Analisi e progettazione di un'interfaccia grafica per la consultazione dei contenuti informativi in una piattaforma web tematica\xspace}
\newcommand{\myDegree}{Doktor-Ingenieur (Dr.-Ing.)\xspace}
\newcommand{\myName}{Nicola Moretto\xspace}
%\newcommand{\myProf}{Put name here\xspace}
%\newcommand{\myOtherProf}{Put name here\xspace}
\newcommand{\mySupervisor}{Prof. Tullio Vardanega\xspace}
\newcommand{\myFaculty}{Facolt\'a di Scienze Matematiche, Fisiche e Naturali\xspace}
\newcommand{\myDepartment}{Dipartimento di Informatica\xspace}
\newcommand{\myUni}{Università di Padova\xspace}
\newcommand{\myLocation}{Padova\xspace}
\newcommand{\myTime}{\today\xspace}
\newcommand{\myVersion}{versione 0.1\xspace}

% ********************************************************************
% Setup, finetuning, and useful commands
% ********************************************************************
\newcounter{dummy} % necessary for correct hyperlinks (to index, bib, etc.)
\newlength{\abcd} % for ab..z string length calculation
\providecommand{\mLyX}{L\kern-.1667em\lower.25em\hbox{Y}\kern-.125emX\@}
\newcommand{\ie}{i.\,e.}
\newcommand{\Ie}{I.\,e.}
\newcommand{\eg}{e.\,g.}
\newcommand{\Eg}{E.\,g.} 
% ****************************************************************************************************


% ****************************************************************************************************
% 3. Loading some handy packages
% ****************************************************************************************************
% ******************************************************************** 
% Packages with options that might require adjustments
% ********************************************************************
\PassOptionsToPackage{square,numbers}{natbib}
 \usepackage{natbib}
\PassOptionsToPackage{fleqn}{amsmath}		% math environments and more by the AMS 
 \usepackage{amsmath}

% ******************************************************************** 
% General useful packages
% ******************************************************************** 
\PassOptionsToPackage{T1}{fontenc} % T2A for cyrillics
	\usepackage{fontenc}                 
\usepackage{xspace} % to get the spacing after macros right  
\usepackage{mparhack} % get marginpar right
\usepackage{fixltx2e} % fixes some LaTeX stuff 
\PassOptionsToPackage{smaller}{acronym}
	\usepackage{acronym} % nice macros for handling all acronyms in the thesis
%\renewcommand*{\acsfont}[1]{\textssc{#1}} % for MinionPro
\renewcommand{\bflabel}[1]{{#1}\hfill} % fix the list of acronyms
% ****************************************************************************************************


% ****************************************************************************************************
% 4. Setup floats: tables, (sub)figures, and captions
% ****************************************************************************************************
\usepackage{tabularx} % better tables
	\setlength{\extrarowheight}{3pt} % increase table row height
\newcommand{\tableheadline}[1]{\multicolumn{1}{c}{\spacedlowsmallcaps{#1}}}
\newcommand{\myfloatalign}{\centering} % to be used with each float for alignment
\usepackage{caption}
\captionsetup{format=hang,font=small}
\usepackage{subfig}  
% ****************************************************************************************************


% ****************************************************************************************************
% 5. Setup code listings
% ****************************************************************************************************
\usepackage{listings} 
%\lstset{emph={trueIndex,root},emphstyle=\color{BlueViolet}}%\underbar} % for special keywords
\lstset{language=[LaTeX]Tex,%C++,
    keywordstyle=\color{RoyalBlue},%\bfseries,
    basicstyle=\small\ttfamily,
    %identifierstyle=\color{NavyBlue},
    commentstyle=\color{Green}\ttfamily,
    stringstyle=\rmfamily,
    numbers=none,%left,%
    numberstyle=\scriptsize,%\tiny
    stepnumber=5,
    numbersep=8pt,
    showstringspaces=false,
    breaklines=true,
    frameround=ftff,
    frame=single,
    belowcaptionskip=.75\baselineskip
    %frame=L
} 
% ****************************************************************************************************    		   


% ****************************************************************************************************
% 6. PDFLaTeX, hyperreferences and citation backreferences
% ****************************************************************************************************
% ********************************************************************
% Using PDFLaTeX
% ********************************************************************
\PassOptionsToPackage{pdftex,hyperfootnotes=false,pdfpagelabels}{hyperref}
	\usepackage{hyperref}  % backref linktocpage pagebackref
\pdfcompresslevel=9
\pdfadjustspacing=1 
\PassOptionsToPackage{pdftex}{graphicx}
	\usepackage{graphicx} 

% ********************************************************************
% Setup the style of the backrefs from the bibliography
% (translate the options to any language you use)
% ********************************************************************
\newcommand{\backrefnotcitedstring}{\relax}%(Not cited.)
\newcommand{\backrefcitedsinglestring}[1]{(Cited on page~#1.)}
\newcommand{\backrefcitedmultistring}[1]{(Cited on pages~#1.)}
\ifthenelse{\boolean{enable-backrefs}}%
{%
		\PassOptionsToPackage{hyperpageref}{backref}
		\usepackage{backref} % to be loaded after hyperref package 
		   \renewcommand{\backreftwosep}{ and~} % separate 2 pages
		   \renewcommand{\backreflastsep}{, and~} % separate last of longer list
		   \renewcommand*{\backref}[1]{}  % disable standard
		   \renewcommand*{\backrefalt}[4]{% detailed backref
		      \ifcase #1 %
		         \backrefnotcitedstring%
		      \or%
		         \backrefcitedsinglestring{#2}%
		      \else%
		         \backrefcitedmultistring{#2}%
		      \fi}%
}{\relax}    

% ********************************************************************
% Hyperreferences
% ********************************************************************
\hypersetup{%
    %draft,	% = no hyperlinking at all (useful in b/w printouts)
    colorlinks=true, linktocpage=true, pdfstartpage=3, pdfstartview=FitV,%
    % uncomment the following line if you want to have black links (e.g., for printing)
    %colorlinks=false, linktocpage=false, pdfborder={0 0 0}, pdfstartpage=3, pdfstartview=FitV,% 
    breaklinks=true, pdfpagemode=UseNone, pageanchor=true, pdfpagemode=UseOutlines,%
    plainpages=false, bookmarksnumbered, bookmarksopen=true, bookmarksopenlevel=1,%
    hypertexnames=true, pdfhighlight=/O,%nesting=true,%frenchlinks,%
    urlcolor=webbrown, linkcolor=RoyalBlue, citecolor=webgreen, %pagecolor=RoyalBlue,%
    %urlcolor=Black, linkcolor=Black, citecolor=Black, %pagecolor=Black,%
    pdftitle={\myTitle},%
    pdfauthor={\textcopyright\ \myName, \myUni, \myFaculty},%
    pdfsubject={},%
    pdfkeywords={},%
    pdfcreator={pdfLaTeX},%
    pdfproducer={LaTeX with hyperref and classicthesis}%
}   

% ********************************************************************
% Setup autoreferences
% ********************************************************************
% There are some issues regarding autorefnames
% http://www.ureader.de/msg/136221647.aspx
% http://www.tex.ac.uk/cgi-bin/texfaq2html?label=latexwords
% you have to redefine the makros for the 
% language you use, e.g., american, ngerman
% (as chosen when loading babel/AtBeginDocument)
% ********************************************************************
\makeatletter
\@ifpackageloaded{babel}%
    {%
       \addto\extrasamerican{%
					\renewcommand*{\figureautorefname}{Figure}%
					\renewcommand*{\tableautorefname}{Table}%
					\renewcommand*{\partautorefname}{Part}%
					\renewcommand*{\chapterautorefname}{Chapter}%
					\renewcommand*{\sectionautorefname}{Section}%
					\renewcommand*{\subsectionautorefname}{Section}%
					\renewcommand*{\subsubsectionautorefname}{Section}% 	
				}%
       \addto\extrasngerman{% 
					\renewcommand*{\paragraphautorefname}{Absatz}%
					\renewcommand*{\subparagraphautorefname}{Unterabsatz}%
					\renewcommand*{\footnoteautorefname}{Fu\"snote}%
					\renewcommand*{\FancyVerbLineautorefname}{Zeile}%
					\renewcommand*{\theoremautorefname}{Theorem}%
					\renewcommand*{\appendixautorefname}{Anhang}%
					\renewcommand*{\equationautorefname}{Gleichung}%        
					\renewcommand*{\itemautorefname}{Punkt}%
				}%	
			% Fix to getting autorefs for subfigures right (thanks to Belinda Vogt for changing the definition)
			\providecommand{\subfigureautorefname}{\figureautorefname}%  			
    }{\relax}
\makeatother


% ****************************************************************************************************
% 7. Last calls before the bar closes
% ****************************************************************************************************
% ********************************************************************
% Development Stuff
% ********************************************************************
\listfiles
%\PassOptionsToPackage{l2tabu,orthodox,abort}{nag}
%	\usepackage{nag}
%\PassOptionsToPackage{warning, all}{onlyamsmath}
%	\usepackage{onlyamsmath}

% ********************************************************************
% Last, but not least...
% ********************************************************************
\usepackage{classicthesis} 
% ****************************************************************************************************


% ****************************************************************************************************
% 8. Further adjustments (experimental)
% ****************************************************************************************************
% ********************************************************************
% Changing the text area
% ********************************************************************
%\linespread{1.05} % a bit more for Palatino
%\areaset[current]{312pt}{761pt} % 686 (factor 2.2) + 33 head + 42 head \the\footskip
%\setlength{\marginparwidth}{7em}%
%\setlength{\marginparsep}{2em}%

% ********************************************************************
% Using different fonts
% ********************************************************************
%\usepackage[oldstylenums]{kpfonts} % oldstyle notextcomp
%\usepackage[osf]{libertine}
%\usepackage{hfoldsty} % Computer Modern with osf
%\usepackage[light,condensed,math]{iwona}
%\renewcommand{\sfdefault}{iwona}
%\usepackage{lmodern} % <-- no osf support :-(
%\usepackage[urw-garamond]{mathdesign} <-- no osf support :-(
% ****************************************************************************************************

\hyphenation{}
\makeindex

\begin{document}
\pagenumbering{arabic}
\pagestyle{scrheadings} %alternativa: headings%
\frenchspacing
\raggedbottom
%\renewcommand*{\bibname}{new name}
%\setbibpreamble{}

%********************************************************************
% Frontmatter
%*******************************************************
%*******************************************************
% Titlepage
%*******************************************************
\begin{titlepage}
	% if you want the titlepage to be centered, uncomment and fine-tune the line below (KOMA classes environment)
	%\begin{addmargin}[-1cm]{-3cm}
    \begin{center}
        \large  

        \hfill

        \vfill

        \begingroup
            \color{Maroon}\spacedallcaps{\myTitle} \\ \bigskip
        \endgroup

        \spacedlowsmallcaps{\myName}

        \vfill

        \includegraphics[width=5cm]{gfx/unipd-logo.png} \\ \medskip

        \mySubtitle \\ \medskip   
        %\myDegree \\
        %\myDepartment \\                            
        %\myFaculty \\
        %\myUni \\ \bigskip

        \myTime\ -- \myVersion

        \vfill                      

    \end{center}  
  %\end{addmargin}       
\end{titlepage}   

\thispagestyle{empty}

\hfill

\vfill

\noindent \textit{\myTitle} \\
\mySubtitle \\
\textcopyright\ \myTime

\medskip

\noindent
\spacedlowsmallcaps{Studente}: \\
\myName (matr. 578258)

\medskip

\noindent\spacedlowsmallcaps{Relatore}: \\
\mySupervisor

\medskip

\noindent\spacedlowsmallcaps{Luogo}: \\
\myLocation

\medskip

\noindent\spacedlowsmallcaps{Data}: \\
\myTime

%*******************************************************
% Dedication
%*******************************************************
\thispagestyle{empty}
%\phantomsection
\refstepcounter{dummy}
\pdfbookmark[1]{Dedica}{Dedication}

\vspace*{3cm}

\begin{flushright}
	\em{Dedicato alla mia famiglia e agli amici di sempre.}   
\end{flushright}

%%*******************************************************
% Abstract
%*******************************************************
%\renewcommand{\abstractname}{Abstract}
\pdfbookmark[1]{Sommario}{Sommario}
\begingroup
\let\clearpage\relax
\let\cleardoublepage\relax
\let\cleardoublepage\relax

\chapter*{Sommario}
	L'attività di stage si è svolta presso l'azienda \textit{Sintesi Sas}, che opera nel settore ICT (\textit{Information and Comunication Technology}) realizzando software ERP e piattaforme Web per aziende, in particolare attive nel settore turistico, e fornendo servizi di consulenza e di formazione di imprenditori nell'ambito del marketing strategico, operativo e del controllo di gestione.
	
	Il prodotto di punta dell'azienda - \textit{Planet Hotel} - costituisce un sistema software per la gestione alberghiera tra i più flessibili, ampi e completi presenti nel panorama italiano, in grado di coprire la maggior parte delle necessità aziendali: oltre alla gestione delle prenotazioni e dei conti, esso offre un insieme di moduli integrati per supportare il controllo di gestione e degli interventi di marketing.

	Si tratta di una realtà imprenditoriale a clientela nazionale con sede unica a Mestre (VE), la cui direzione e amministrazione è affidata al solo fondatore, che ha assunto il ruolo di tutor esterno e referente aziendale per l'intera durata dello stage.

	Le attività svolte si inseriscono nell'ambito di un progetto esterno rispetto al business dell'azienda, finalizzato alla realizzazione di una piattaforma web tematica per la condivisione di informazioni e la vendita diretta di prodotti alla clientela e affidato ad un team costituito da differenti figure professionali (sociologi, informatici, ingegneri, \ldots).
	
	\section*{Contenuti}
	\label{ch:tesi:intro:contenuti}
	Il presente documento costituisce una relazione dettagliata in merito all'attività di stage svolta dallo studente Nicola Moretto presso l'azienda \textit{Sintesi Sas}. I contenuti sono organizzati nei seguenti capitoli:
	\begin{description}
	  \item[\nameref{ch:tesi:progetto}] \hfill \\
	  Il primo capitolo illustra le strategie dell'azienda e gli obiettivi, i requisiti e i vincoli del progetto in cui si inseriscono le attività di stage.
	  \item[\nameref{ch:tesi:stage}] \hfill \\ 
		Il secondo capitolo illustra gli obiettivi, i requisiti e l'organizzazione (piano e norme di lavoro) delle attività di stage. A seguire vengono presentate le scelte più rilevanti effettuate e i risultati conseguiti.
	  \item[\nameref{ch:tesi:conclusioni}] \hfill \\
		Il terzo capitolo presenta un'analisi critica a posteriori dell'attività di stage: raggiungimento degli obiettivi prefissati, competenze professionali acquisite, \ldots\ .
	\end{description}

	\section*{Convenzioni tipografiche}
	\label{ch:tesi:intro:convenzioni}
	Al fine di agevolare la consultazione del documento sono state adottate alcune convenzioni tipografiche illustrate di seguito.

	\paragraph{Glossario} Gli acronimi, le abbreviazioni, i nomi propri e i termini specialistici contenuti nel presente documento sono illustrati nel \textit{\nameref{ch:tesi:appendice:glossario}}, consultabile in appendice, al fine di agevolare la lettura e la comprensione degli argomenti trattati. La prima occorrenza di ciascun termine o espressione presente nel glossario appare \underline{sottolineata}.

	\paragraph{Terminologia} I termini propri o di provenienza straniera divenuti di uso corrente nella lingua italiana sono evidenziati in \textit{corsivo}, mentre la prima occorrenza di parole o espressioni che assumono particolare significato nel presente contesto è riportata in \textsc{maiuscoletto}.

	\paragraph{Codice e formule} I nomi di tabelle, classi, package, \ldots\ impiegano uno stile di carattere \textsf{sans serif}, mentre i frammenti di codice impiegano un carattere a \texttt{spaziatura fissa}.
\endgroup		

\vfill

%%*******************************************************
% Acknowledgments
%*******************************************************
\pdfbookmark[1]{Acknowledgments}{acknowledgments}

\begin{flushright}{\slshape    
    We have seen that computer programming is an art, \\ 
    because it applies accumulated knowledge to the world, \\ 
    because it requires skill and ingenuity, and especially \\
    because it produces objects of beauty.} \\ \medskip
    --- \defcitealias{knuth:1974}{Donald E. Knuth}\citetalias{knuth:1974} \citep{knuth:1974}
\end{flushright}



\bigskip

\begingroup
\let\clearpage\relax
\let\cleardoublepage\relax
\let\cleardoublepage\relax
\chapter*{Ringraziamenti}
Put your acknowledgments here.

Many thanks to everybody who already sent me a postcard!

Regarding the typography and other help, many thanks go to Marco 
Kuhlmann, Philipp Lehman, Lothar Schlesier, Jim Young, Lorenzo 
Pantieri and Enrico Gregorio\footnote{Members of GuIT (Gruppo 
Italiano Utilizzatori di \TeX\ e \LaTeX )}, J\"org Sommer, 
Joachim K\"ostler, Daniel Gottschlag, Denis Aydin, Paride 
Legovini, Steffen Prochnow, Nicolas Repp, Hinrich Harms, 
 Roland Winkler,  
 and the whole \LaTeX-community for support, ideas and 
 some great software.

\bigskip

\noindent\emph{Regarding \mLyX}: The \mLyX\ port was intially done by 
\emph{Nicholas Mariette} in March 2009 and continued by 
\emph{Ivo Pletikosi\'c} in 2011. Thank you very much for your 
work and the contributions to the original style.


\endgroup




%*******************************************************
% Table of Contents
%*******************************************************
%\phantomsection
\refstepcounter{dummy}
\pdfbookmark[1]{\contentsname}{tableofcontents}
\setcounter{tocdepth}{2} % <-- 2 includes up to subsections in the ToC
\setcounter{secnumdepth}{2} % <-- 3 numbers up to subsubsections
\manualmark
\markboth{\spacedlowsmallcaps{\contentsname}}{\spacedlowsmallcaps{\contentsname}}
\tableofcontents 
\automark[section]{chapter}
\renewcommand{\chaptermark}[1]{\markboth{\spacedlowsmallcaps{#1}}{\spacedlowsmallcaps{#1}}}
\renewcommand{\sectionmark}[1]{\markright{\thesection\enspace\spacedlowsmallcaps{#1}}}
%*******************************************************
% List of Figures and of the Tables
%*******************************************************
\clearpage

\begingroup 
    \let\clearpage\relax
    \let\cleardoublepage\relax
    \let\cleardoublepage\relax
    %*******************************************************
    % List of Figures
    %*******************************************************    
    %\phantomsection 
    \refstepcounter{dummy}
    %\addcontentsline{toc}{chapter}{\listfigurename}
    \pdfbookmark[1]{\listfigurename}{lof}
    \listoffigures

    \vspace*{8ex}

    %*******************************************************
    % List of Tables
    %*******************************************************
    %\phantomsection 
    \refstepcounter{dummy}
    %\addcontentsline{toc}{chapter}{\listtablename}
    \pdfbookmark[1]{\listtablename}{lot}
    \listoftables
        
    \vspace*{8ex}
%   \newpage
    
    %*******************************************************
    % List of Listings
    %*******************************************************      
	  %\phantomsection 
    %\refstepcounter{dummy}
    %\addcontentsline{toc}{chapter}{\lstlistlistingname}
    %\pdfbookmark[1]{\lstlistlistingname}{lol}
    %\lstlistoflistings 

    %\vspace*{8ex}
       
    %*******************************************************
    % Acronyms
    %*******************************************************
    %\phantomsection 
    %\refstepcounter{dummy}
    %\pdfbookmark[1]{Abbreviazioni e acronimi}{acronyms}
    %\markboth{\spacedlowsmallcaps{Acronyms}}{\spacedlowsmallcaps{Acronyms}}
    %\chapter*{Abbreviazioni e acronimi}
    %\begin{acronym}
    %\end{acronym}                     
\endgroup

\cleardoublepage


%********************************************************************
% Parte 1
%*******************************************************
% use \cleardoublepage here to avoid problems with pdfbookmark
\ctparttext{Presentazione dei criteri di classificazione addizionali per la catalogazione ed il reperimento dei contenuti informativi pubblicati dagli utenti.}
\part{Sistema di classificazione}
%************************************************
\chapter{Presentazione}\label{ch:presentazione}
%************************************************
Primo paragrafo scritto con \TeX{}Live.
%*****************************************



%*****************************************
\chapter{Relazioni}\label{ch:contenuti:relazioni}
%*****************************************
\ldots
%*****************************************

%************************************************
\chapter{Dizionario}
\label{ch:intro:dizionario}
%************************************************
%*****************************************


%********************************************************************
% Parte 2
%*******************************************************
\ctparttext{Presentazione dell'interfaccia grafica per la visualizzazione, la navigazione e la consultazione dei contenuti informativi.}
\part{Interfaccia grafica}
%%************************************************
\chapter{Eventi}\label{ch:eventi}
%************************************************
Scrivere qui...
%*****************************************

%%************************************************
\chapter{Argomenti}\label{ch:argomenti}
%************************************************
Tassonomia
%*****************************************


% ********************************************************************
% Backmatter
%*******************************************************
\appendix
\part{Appendice}
	%********************************************************************
% Appendix
%*******************************************************
% If problems with the headers: get headings in appendix etc. right
%\markboth{\spacedlowsmallcaps{Appendix}}{\spacedlowsmallcaps{Appendix}}
\chapter{Glossario}
\label{ch:tesi:appendice:glossario}

\section*{B}
\paragraph{Bitbucket - \url{https://bitbucket.org/}} \hfill \\
Piattaforma web per la gestione delle attività di progetto con supporto a strumenti di controllo di versione distribuito.

\section*{C}
\paragraph{CamelCase} \hfill \\
Convenzione per la scrittura di espressioni composte unendo le parole tra loro e mantenendo ciascuna iniziale in maiuscolo.

\section*{G}
\paragraph{gedit - \url{http://projects.gnome.org/gedit/}} \hfill \\
Editor di testo ufficiale dell'ambiente desktop GNOME.
%\paragraph{Gruppi di acquisto} \hfill \\
%Insieme (stabile o provvisorio) di consumatori che acquista mediante ordine collettivo un consistente numero di beni direttamente dal produttore, spesso per conseguire prezzi vantaggiosi o ammortizzare eventuali spese accessorie.

\section*{L}
\paragraph{LaTeXila - \url{http://projects.gnome.org/latexila/}} \hfill \\
Editor LaTex integrato per l'ambiente desktop GNOME.
\paragraph{LibreOffice Calc - \url{http://www.libreoffice.org/}} \hfill \\
Applicazione per fogli di calcolo della suite di produttività \textit{LibreOffice}.

\section*{M}
\paragraph{Mercurial - \url{http://mercurial.selenic.com/}} \hfill \\
Strumento multi piattaforma, gratuito ed open source per il controllo di versione distribuito.
\paragraph{MySQL Workbench - \url{http://www.mysql.it/products/workbench/}} \hfill \\
Applicazione multi piattaforma, gratuita ed open source per la progettazione, lo sviluppo e l'amministrazione di database MySQL. 

\section*{P}
\paragraph{PDF (Portable Document Format)} \hfill \\
Formato di file per la rappresentazione di documenti in maniera indipendente dalla piattaforma hardware e software.
\paragraph{Pencil - \url{http://pencil.evolus.vn/}} \hfill \\
Applicazione multi piattaforma, gratuita ed open source per la realizzazione di prototipi di interfacce grafiche.
\paragraph{ProjectLibre - \url{http://www.projectlibre.org/}} \hfill \\
Applicazione multi piattaforma, gratuita ed open source per il \textit{project management}, che consente di realizzare diagrammi di Gantt e di PERT, gestire le risorse allocate e le attività pianificate, \ldots\ .

\section*{R}
\paragraph{Rete sociale} \hfill \\
Insieme di persone, aventi interessi in comune e inclini a collaborare e condividere idee o informazioni, e di relazioni di tipo esperienziale definite tra tali soggetti.

\section*{U}
\paragraph{Ubuntu - \url{http://www.ubuntu.com/}} \hfill \\
Distribuzione Linux gratuita derivata da Debian.
\paragraph{UML - \url{http://www.uml.org/}} \hfill \\
Standard internazionale per un linguaggio di modellazione, che definisce un insieme di notazioni grafiche per la rappresentazione visiva di sistemi.
\paragraph{UMLet - \url{http://www.umlet.com/}} \hfill \\
Applicazione multi piattaforma, gratuita ed open source per la realizzazione di diagrammi UML.

\cleardoublepage%********************************************************************
% Bibliography
%*******************************************************
% work-around to have small caps also here in the headline
\manualmark
\markboth{\spacedlowsmallcaps{\bibname}}{\spacedlowsmallcaps{\bibname}} % work-around to have small caps also
%\phantomsection 
\refstepcounter{dummy}
\addtocontents{toc}{\protect\vspace{\beforebibskip}} % to have the bib a bit from the rest in the toc
\addcontentsline{toc}{chapter}{\tocEntry{\bibname}}
\bibliographystyle{plainnat}
\label{app:bibliography} 
\bibliography{Bibliography}
%\cleardoublepage\printindex
\end{document}
